\section{流量峰值、Multi-Cloud 与优雅降级}

大型社区平台的峰值往往由突发新闻、体育、政治事件或外部链接引发，时间和地域很难完全预测。基础设施目标不是让所有功能在任何峰值下保持同等质量，而是保护核心读写路径、身份和安全能力，再按业务价值逐步降级昂贵功能。Reddit 的经验可以概括为 absorb、divert、degrade、recover 四步。

\subsection{四步韧性路径}

第一步用缓存、队列和弹性余量吸收普通峰值；第二步通过多区域或多云流量调度绕开局部故障；第三步在容量不足时进入只读或关闭高成本路径；第四步恢复写入、回放队列并核对状态。\term{Graceful degradation} 即优雅降级，指系统在资源不足或部分故障时保留最重要能力，而不是整体崩溃。

\begin{figure}[H]
\centering
\includegraphics[width=0.96\textwidth]{images/09-traffic-resilience.png}
\caption{社区平台的流量韧性路径：吸收、切流、降级与恢复。}
\figsource{依据访谈字幕 40:00--43:20 重绘，`images/09-traffic-resilience.png`。}
\end{figure}

\begin{knowledgebox}{读图：降级顺序必须提前编码}
事故时再讨论哪些功能重要，决策会被压力和局部信息支配。系统应预先定义核心读路径、发帖与评论写路径、登录、安全审核、搜索、推荐和通知的优先级。只读模式能保护社区内容可见性，却也可能阻止版主处理正在扩散的伤害，因此安全工具不应和普通写功能一起被粗暴关闭。
\end{knowledgebox}

\subsection{Multi-cloud 不是复制一套机器}

多云能降低单一供应商故障和议价风险，也会增加数据复制、身份、网络、可观测性和演练成本。只有当团队能够在可接受时间内切流、确认数据一致性并让关键服务在第二环境中独立运行，多云才是恢复能力。把少量资源部署到另一家云却从未演练，更多是采购多样性而非故障切换。

\begin{importantbox}{课堂提示：恢复目标要以用户动作表达}
``第二云可用'' 不够具体。更好的目标是：用户能读取热门帖子，版主能隐藏危险内容，登录状态不会大面积失效，新评论可以排队并在恢复后写入。用户动作能暴露依赖，也能指导演练是否真正覆盖核心价值。
\end{importantbox}

\begin{table}[H]
\centering
\begin{tabular}{L{0.19\textwidth}L{0.31\textwidth}L{0.4\textwidth}}
\toprule
能力 & 失败时的用户影响 & 降级策略 \\
\midrule
内容读取 & 平台几乎失去用途 & 多层缓存、静态快照、跨区只读 \\
发帖评论 & 实时参与中断 & 排队、限流、延迟写入 \\
版主安全工具 & 伤害无法及时控制 & 保留专用容量与最小审核接口 \\
推荐排序 & 个性化下降 & 回退到社区热门或时间排序 \\
搜索/AI 摘要 & 信息发现变慢 & 关闭摘要，保留关键词检索 \\
通知与分析 & 延迟可接受 & 批处理和恢复后回放 \\
\bottomrule
\end{tabular}
\caption{按用户价值设计降级，而不是按微服务所有权设计降级。}
\end{table}

\begin{warningbox}{恢复后的数据修复同样重要}
流量切换成功不代表事故结束。队列可能重复消费，计数器可能漂移，投票顺序和评论树可能出现短暂不一致，安全事件也可能在降级窗口积压。恢复流程需要 reconciliation、去重、审计和用户补偿，不能只看 HTTP 成功率回升。
\end{warningbox}

\subsection{本章小结}

平台韧性来自预留容量、可演练的切流、按用户价值排序的优雅降级和恢复后数据核对。Multi-cloud 只有在应用、数据和身份都能真实切换时才提供韧性。接下来讨论 AI 如何改变流量入口与接口形态，同时哪些社区需求不会被摘要模型替代。

